\Needspace{8\baselineskip}
\Article{282}
Das Auffinden der Darstellungen einer gegebenen binären Form, deren
Determinante nicht $=0$ ist\HistoricalNote{*)}{Diesen Fall, der eine etwas andere Methode erfordert, übergehen wir hier der Kürze halber.},
durch eine gegebene ternäre Form beruht auf den folgenden Bemerkungen:

I.\ Aus jeder eigentlichen Darstellung der binären Form $(p,q,r)=\varphi$ mit der
Determinante $D$ durch die ternäre Form $f$ mit der Determinante
$\Delta$ lassen sich ganze Zahlen $B,B'$ ableiten, für die
\[
BB\equiv\Delta p,\qquad BB'\equiv-\Delta q,\qquad
B'B'\equiv\Delta r\ (\mathrm{mod.}\ D)
\]
gilt, \textit{d.h.}, ein Wert des Ausdrucks $\surd\Delta(p,\allowbreak -q,\allowbreak r)(\mathrm{mod.}\ D)$.
Eine eigentliche Darstellung der Form $\varphi$ durch $f$ sei
\[
x=\alpha t+\beta u,\qquad x'=\alpha't+\beta'u,\qquad x''=\alpha''t+\beta''u
\]
(wobei $x,x',x''$; $t,u$ die Unbestimmten der Formen $f,\varphi$ bezeichnen).
Man wähle ganze Zahlen $\gamma,\gamma',\gamma''$ so, dass
\[
(\alpha'\beta''-\alpha''\beta')\gamma+
(\alpha''\beta-\alpha\beta'')\gamma'+
(\alpha\beta'-\alpha'\beta)\gamma''
\]
$=k$ entweder $=+1$ oder $=-1$ ist, und die Form $f$ gehe
durch die Substitution
\[
\begin{array}{ccc}
\alpha,&\beta,&\gamma\\
\alpha',&\beta',&\gamma'\\
\alpha'',&\beta'',&\gamma''
\end{array}
\]
in die Form
$\left(\begin{smallmatrix}a,a',a''\\b,b',b''\end{smallmatrix}\right)=g$
über, der die Form
$\left(\begin{smallmatrix}A,A',A''\\B,B',B''\end{smallmatrix}\right)=G$
adjungiert sei. Dann sind offenbar $a=p$, $b''=q$, $a'=r$, $A''=D$,
und $\Delta$ ist die Determinante der Form $g$; also gilt
\[
BB=\Delta p+A'D,\qquad BB'=-\Delta q+B''D,\qquad B'B'=\Delta r+AD
\]
So wird \textit{z.B.}\ die Form $19tt+6tu+41uu$ durch $xx+x'x'+x''x''$ dargestellt,
indem man $x=3t+5u$, $x'=3t-4u$, $x''=t$ setzt. Setzt man dann $\gamma=-1$,
$\gamma'=1$, $\gamma''=0$,
\SourceLeaf{329}
so erhält man $B=-171$, $B'=27$, also den Wert $(-171,27)$ des Ausdrucks
$\surd-1(19,\allowbreak -3,\allowbreak 41)(\mathrm{mod.}\ 770)$.

Daraus folgt nun, dass, wenn $\Delta(p,-q,r)$ kein quadratischer Rest von $D$
ist, $\varphi$ durch keine ternäre Form mit der Determinante $\Delta$
eigentlich dargestellt werden kann. Sind also $\Delta,D$ teilerfremd,
so muss $\Delta$ eine charakteristische Zahl der Form $\varphi$ sein.

II.\ Da man $\gamma,\gamma',\gamma''$ auf unendlich viele verschiedene Arten bestimmen
kann, ergeben sich daraus auch immer wieder andere Werte von $B,B'$.
Sehen wir uns ihren Zusammenhang an. Nehmen wir an, auch $\delta,\delta',\delta''$ seien
so beschaffen, dass
\[
(\alpha'\beta''-\alpha''\beta')\delta+
(\alpha''\beta-\alpha\beta'')\delta'+
(\alpha\beta'-\alpha'\beta)\delta''=\mathfrak{k}
\]
entweder $=+1$ oder $=-1$ ist, und die Form $f$ gehe durch
die Substitution
\[
\begin{array}{ccc}
\alpha,&\beta,&\delta\\
\alpha',&\beta',&\delta'\\
\alpha'',&\beta'',&\delta''
\end{array}
\]
in
$\left(\begin{smallmatrix}\mathfrak{a},\mathfrak{a}',\mathfrak{a}''\\\mathfrak{b},\mathfrak{b}',\mathfrak{b}''\end{smallmatrix}\right)=\mathfrak{g}$
über, der die Form
$\left(\begin{smallmatrix}\mathfrak{A},\mathfrak{A}',\mathfrak{A}''\\\mathfrak{B},\mathfrak{B}',\mathfrak{B}''\end{smallmatrix}\right)=\mathfrak{G}$
adjungiert sei. Dann sind $g,\mathfrak{g}$ äquivalent, also auch $G$ und $\mathfrak{G}$.
Wendet man die Regeln aus Art.\ 269, 270 an\HistoricalNote{*)}{Man leite aus der Transformation der Form $f$ in $g$ die Transformation der Form $g$ in $f$ ab; aus dieser und der Transformation der Form $f$ in $\mathfrak{g}$ die Transformation der Form $g$ in $\mathfrak{g}$; schließlich aus dieser durch Transposition die Transformation der Form $\mathfrak{G}$ in $G$.},
so findet man, wenn man
\[
\begin{aligned}
(\beta'\gamma''-\beta''\gamma')\delta+
(\beta''\gamma-\beta\gamma'')\delta'+
(\beta\gamma'-\beta'\gamma)\delta''&=\zeta\\
(\gamma'\alpha''-\gamma''\alpha')\delta+
(\gamma''\alpha-\gamma\alpha'')\delta'+
(\gamma\alpha'-\gamma'\alpha)\delta''&=\eta
\end{aligned}
\]
setzt, dass die Form $\mathfrak{G}$ in $G$ übergeht, und zwar durch die Substitution
\[
\begin{array}{ccc}
k,&0,&0\\
0,&k,&0\\
\zeta,&\eta,&\mathfrak{k}
\end{array}
\]
Daraus folgt
\[
B=\eta\mathfrak{k}D+\mathfrak{k}k\mathfrak{B},\qquad
B'=\zeta\mathfrak{k}D+\mathfrak{k}k\mathfrak{B}'
\]
und also wegen $\mathfrak{k}k=\pm1$ entweder $B\equiv\mathfrak{B}$, $B'\equiv\mathfrak{B}'$ oder $B\equiv-\mathfrak{B}$, $B'\equiv-\mathfrak{B}'$
$(\mathrm{mod.}\ D)$. Im ersten Fall nennen wir die Werte $(B,B')$, $(\mathfrak{B},\mathfrak{B}')$ äquivalent,
im zweiten entgegengesetzt. Von der Darstellung der Form $\varphi$ sagen
wir, dass sie zu jedem Wert des
\SourceLeaf{330}
Ausdrucks $\surd\Delta(p,\allowbreak -q,\allowbreak r)(\mathrm{mod.}\ D)$ \textit{gehört}, der sich aus ihr nach der Methode in I
ableiten lässt. Alle Werte, zu denen dieselbe Darstellung gehört, sind
also entweder äquivalent oder entgegengesetzt.

III.\ Wenn umgekehrt, wie oben in I, die Darstellung der Form $\varphi$
durch $f$ mittels $x=\alpha t+\beta u$ usw.\ zu dem Wert $(B,B')$ gehört, der
aus ihr mithilfe der Transformation
\[
\begin{array}{ccc}
\alpha,&\beta,&\gamma\\
\alpha',&\beta',&\gamma'\\
\alpha'',&\beta'',&\gamma''
\end{array}
\]
abgeleitet wird, so gehört dieselbe auch zu jedem anderen Wert $(\mathfrak{B},\mathfrak{B}')$,
der dem ersten äquivalent oder entgegengesetzt ist;
\textit{d.h.}, anstelle von $\gamma,\gamma',\gamma''$ kann man andere ganze Zahlen
$\delta,\delta',\delta''$ wählen, für die die Gleichung $(\Omega)$, nämlich
\[
(\alpha'\beta''-\alpha''\beta')\delta+
(\alpha''\beta-\alpha\beta'')\delta'+
(\alpha\beta'-\alpha'\beta)\delta''=\pm1
\]
gilt und die so beschaffen sind, dass die Koeffizienten mit den Nummern
$4$ und $5$ in der Form, die derjenigen adjungiert ist, in die
$f$ durch die Substitution $(S)$
\[
\begin{array}{ccc}
\alpha,&\beta,&\delta\\
\alpha',&\beta',&\delta'\\
\alpha'',&\beta'',&\delta''
\end{array}
\]
übergeht, jeweils $=\mathfrak{B},\mathfrak{B}'$ werden. Man setze nämlich
\[
\pm B=\mathfrak{B}+\eta D,\qquad \pm B'=\mathfrak{B}'+\zeta D
\]
(wobei hier und im Folgenden die oberen oder unteren Vorzeichen zu nehmen
sind, je nachdem die Werte $(B.\ B')$, $(\mathfrak{B},\mathfrak{B}')$ äquivalent oder entgegengesetzt
sind). Dann sind $\zeta,\eta$ ganze Zahlen. Die Form $g$ gehe durch
die Substitution
\[
\begin{array}{ccc}
1,&0,&\zeta\\
0,&1,&\eta\\
0,&0,&\pm1
\end{array}
\]
in die Form $\mathfrak{g}$ über. Man sieht leicht, dass ihre Determinante
$\Delta$ ist und dass in der adjungierten Form die Koeffizienten mit den
Nummern $4$ und $5$ jeweils $=\mathfrak{B},\mathfrak{B}'$ werden. Setzt man aber
\[
\alpha\zeta+\beta\eta\pm\gamma=\delta,\quad
\alpha'\zeta+\beta'\eta\pm\gamma'=\delta',\quad
\alpha''\zeta+\beta''\eta\pm\gamma''=\delta''
\]
\SourceLeaf{331}
so sieht man ohne Schwierigkeit, dass $f$ durch die Substitution
$(S)$ in $\mathfrak{g}$ übergeht und dass die Gleichung $(\Omega)$ erfüllt
ist. \textit{Q. E. D.}

